\clearpage
\textbf{5. Auswahlregeln}\\
Aus Betrachtung der Symmetrieeigenschaften des Dipol-Übergangsmatrixelements (QM)

$$
\begin{aligned}
	\Delta \Lambda &= 0, \pm 1 \\
	\Delta S &= 0 \\
	\Delta \Omega &= 0, \pm 1 
\end{aligned} 
\qquad
\begin{aligned}
+ &\to + \\
- &\to - \\
\mathrm{u} &\to \mathrm{g}
\end{aligned}
$$

\textbf{6. Potentialkurven}\\
Die Elektronenanordnung beeinflusst die Bindungsverhältnisse zwischen den Atomkernen. Daraus folgt, dass zu jedem Elektronenzustand des Moleküls eine Potentialkurve, also ein Morsepotential, gehört.

\begin{figure}[H]
    \centering
    \incfig{vl12_abbildung_1}
    \label{fig:vl12_abbildung_1}
\end{figure}

\section{Elektronenspektren von Molekülen}
$$
\Delta T = \Delta T_{\mathrm{e}} + \Delta G(\nu) + \Delta F(J)
$$ 
$$
\begin{aligned}
	\Delta T_{\mathrm{e}} &\gg \Delta G(\nu) &&> \Delta F(J) \\
	\unit{\electronvolt}& \phantom{ \gg}\si{10^{-1} \electronvolt} - \si{10^{-2} \electronvolt} && \quad \si{10^{-3} \electronvolt}
\end{aligned}
$$ 
\begin{itemize}
	\item Elektronenübergang holt das Rotationsschwingungsspektrum in den sichtbaren Spektralbereich 
	\item auch für homonukleare Moleküle (da sich beim Elektronenübergang auch das Dipolmoment ändern kann)
\end{itemize}

\textbf{a. Elektronenübergang und Schwingungsänderung}

\begin{figure}[H]
    \centering
    \incfig{vl12_abbildung_2}
    \label{fig:vl12_abbildung_2}
\end{figure}
Die alte Auswahlregel
$$
\Delta \nu = \pm 1, \pm 2, \pm 3
$$ 
ist \textbf{nicht} mehr gültig! Das \textbf{Franck-Condon Prinzip} macht jetzt Aussagen:\\
Klassisch erfolgt ein Elektronenübergang schnell gegenüber der Schwingung, da die Elektronenmasse sehr viel geringer ist als die Kernmasse.
Daher muss der Abstand der Kerne der gleiche sein, da sich der Schwingungszustand nicht ändert.
\begin{enumerate}
	\item $r = \text{const}$ $\implies$ Übergang muss in der Abbildung \textbf{senkrecht} sein
	\item Von Umkehrpunkt zu Umkehrpunkt, da die Aufenthaltswahrscheinlichkeit an den Umkehrpunkten am größten ist. Die Ausnahme hierbei bildet der Grundzustand.
\end{enumerate}

\begin{figure}[H]
    \centering
    \incfig{vl12_abbildung_3}
    \label{fig:vl12_abbildung_3}
\end{figure}

Quantentheorie:
$$
W \sim \int \chi_{\nu'} (r) \xi_{\nu''}(r) \mathrm{d}V
$$ 
Überlappintegral der Kernschwingungsfunktion $\chi_{\nu'}$ und $\chi_{\nu''}$, bei der Auslenkung $r$, integriert über das Molekülvolumen.\\
Wegen der endlichen Breite der Wahrscheinlichkeitsbereiche gibt es nicht nur einen Übergang sondern man erhält mit unterschiedlichen Wahrscheinlichkeiten Übergänge zu benachbarten Schwingungsniveaus.\\

\textbf{Aussehen der Spektren (Absorbtion)}\\
\textbf{Fall 1}: $r_{\mathrm{e}}' \simeq r_{\mathrm{e}}''$ \\
\begin{figure}[H]
    \centering
    \incfig{vl12_abbildung_4}
    \label{fig:vl12_abbildung_4}
\end{figure}


\textbf{Fall 2}: $r_{\mathrm{e}}'=r_{\mathrm{e}}''$\\
\begin{figure}[H]
    \centering
    \incfig{vl12_abbildung_5}
    \label{fig:vl12_abbildung_5}
\end{figure}

\textbf{c. Emissionspektren}\\
Im Prinzip die gleichen Bandenspektren wie in Absorbtion.\\
Bsp. $r_{\mathrm{e}}' \neq r_{\mathrm{e}}''$ \\

\begin{figure}[H]
    \centering
    \incfig{vl12_abbildung_6}
    \label{fig:vl12_abbildung_6}
\end{figure}

\begin{figure}[H]
    \centering
    \def\svgwidth{.35\columnwidth}
    \subimport{./figures}{vl12_abbildung_7.pdf_tex}
    \label{fig:vl12_abbildung_7}
\end{figure}

Strahlungslose Übergänge durch Stöße! \\
Konkurenzprozess zu strahlenden Übergängen $W_{\mathrm{opt}} \sim (h \nu)^3$
im Infrarot bereich klein $\implies$ Relaxation strahlungslos in den Grundzustand\\

Anwendung: Farbstofflaser (4 Niveaus)

\textbf{d. Rotationsstruktur}
\begin{figure}[H]
    \centering
    \def\svgwidth{.31\columnwidth}
    \subimport{./figures}{vl12_abbildung_9.pdf_tex}
    \label{fig:vl12_abbildung_7}
\end{figure}

$$
\overline{\nu} = T' - T'' = T' ^{\text{el}} G_{\nu'}' + F'_{\nu', J'} - \left( T''^{\text{el}} + G_{\nu''}'' + F_{\nu'', J''}'' \right) 
$$ 
Vernachlässigung des Zentrifugalterms $(D_{\nu}J^2(J+1)^2)$

$$
\begin{aligned}
	F_{\nu', J'}' &= B_{\nu'}J'(J'+1)\\
	F_{\nu'', J''}'' &= B_{\nu''}J'' (J''+1)
\end{aligned}
$$ 
